\begin{figure}[tb]
\centering
\begin{tikzpicture}
\begin{axis}[axis main,width=.8\linewidth,height=6.2cm,
 xlabel={$\nu=U^2/P_{\mathrm{Cu}}$},ylabel={$\tau=M/P_{\mathrm{Cu}}$},
 xmin=0,xmax=5.2,ymin=-2.1,ymax=2.8]
 \addplot[derived quantity,strong line,derived light region,
 domain=0:360,samples=161]
 ({2.4+1.75*cos(x)},{.25+1.35*cos(x)*.25+1.35*sin(x)*.97});
 \addplot[support line] coordinates {(.5,2.55) (4.8,1.25)};
 \addplot[voltage curve] coordinates {(0,0) (4.9,2.45)};
 \addplot[only marks,mark=*,estimated quantity] coordinates {(1.55,2.23)};
 \node[annotation] at (axis cs:3.9,2.0){voltage-feasibility line};
 \node[annotation] at (axis cs:1.55,-1.45){compact convex direction image};
\end{axis}
\end{tikzpicture}
\caption{Every point represents a current direction. A rotating support line
finds the upper boundary; the requested-torque voltage inequality is a line
through the origin. Their contact selects the optimal direction.}
\label{fig:performance-space}
\end{figure}
